\documentclass[12pt]{ximera}
\typeout{Start loading xmPreamble.tex}%

% Add here extra macro's that are loaded automatically by all documents of claas 'ximera' or 'xourse' in this repo

%%
%%  Example:
%%
% \newcommand{\R}{\mathbb{R}}
\usepackage{tikz}
\usetikzlibrary{positioning, arrows.meta}
\usepackage{multicol}
\usepackage{enumitem}
\usepackage{caption}
\usepackage{amsmath}
\usepackage{amsthm}
\usepackage{fontenc}

%% ---------------------------------------------------------------
%% Shared worksheet macros.
%%
%% These came from the Summer 2026 worksheet set, where each worksheet carried its
%% own copy. They have to live here instead: a xourse inlines every activity into a
%% single document, so duplicate \newcommand definitions across activities collide,
%% and the xourse gobbles \input so a per-topic macro file never loads.
%%
%% They use \providecommand, NOT \newcommand. ximera.cls loads this file automatically,
%% and every activity also carries an explicit \typeout{Start loading xmPreamble.tex}%

% Add here extra macro's that are loaded automatically by all documents of claas 'ximera' or 'xourse' in this repo

%%
%%  Example:
%%
% \newcommand{\R}{\mathbb{R}}
\usepackage{tikz}
\usetikzlibrary{positioning, arrows.meta}
\usepackage{multicol}
\usepackage{enumitem}
\usepackage{caption}
\usepackage{amsmath}
\usepackage{amsthm}
\usepackage{fontenc}

%% ---------------------------------------------------------------
%% Shared worksheet macros.
%%
%% These came from the Summer 2026 worksheet set, where each worksheet carried its
%% own copy. They have to live here instead: a xourse inlines every activity into a
%% single document, so duplicate \newcommand definitions across activities collide,
%% and the xourse gobbles \input so a per-topic macro file never loads.
%%
%% They use \providecommand, NOT \newcommand. ximera.cls loads this file automatically,
%% and every activity also carries an explicit \typeout{Start loading xmPreamble.tex}%

% Add here extra macro's that are loaded automatically by all documents of claas 'ximera' or 'xourse' in this repo

%%
%%  Example:
%%
% \newcommand{\R}{\mathbb{R}}
\usepackage{tikz}
\usetikzlibrary{positioning, arrows.meta}
\usepackage{multicol}
\usepackage{enumitem}
\usepackage{caption}
\usepackage{amsmath}
\usepackage{amsthm}
\usepackage{fontenc}

%% ---------------------------------------------------------------
%% Shared worksheet macros.
%%
%% These came from the Summer 2026 worksheet set, where each worksheet carried its
%% own copy. They have to live here instead: a xourse inlines every activity into a
%% single document, so duplicate \newcommand definitions across activities collide,
%% and the xourse gobbles \input so a per-topic macro file never loads.
%%
%% They use \providecommand, NOT \newcommand. ximera.cls loads this file automatically,
%% and every activity also carries an explicit \input{../xmPreamble}, so this file is
%% read twice. That was harmless while it held only \usepackage lines (LaTeX skips a
%% package it has already loaded) but a second \newcommand is a hard error.
%%
%%   \blank[width]        fill-in-the-blank rule (default 2.5cm)
%%   \showwork            "show and explain your work" reminder
%%   \showworkline        ... plus which schedule line was used
%%   \showworkyear        ... plus which year's schedule was used
%%   \schedhead{...}      centred bold caption above a tiered-rate table
%%   \samesquares[scale]{left}{right}   two equal-area squares, labelled
%%   \samecubes[scale]{left}{right}     two equal-volume cubes, labelled
%% ---------------------------------------------------------------

\providecommand{\blank}[1][2.5cm]{\underline{\hspace{#1}}}
\providecommand{\showwork}{\textbf{REMEMBER TO SHOW AND EXPLAIN YOUR WORK.}}
\providecommand{\showworkline}{\textbf{REMEMBER TO SHOW AND EXPLAIN YOUR WORK.
  Explanation should include how and why you know which line of the
  schedule was used to calculate this amount.}}
\providecommand{\showworkyear}{\begin{sloppypar}\textbf{REMEMBER TO SHOW AND
  EXPLAIN YOUR WORK. Explanation should include how and why you know which
  line of the year's tax schedule was used to calculate this tax
  amount.}\end{sloppypar}}
\providecommand{\schedhead}[1]{\begin{center}\textbf{#1}\end{center}}

\providecommand{\samesquares}[3][1]{%
\begin{center}
{\footnotesize These squares are the \textbf{same size}.}

\vspace{1.5mm}
\begin{tikzpicture}[scale=#1,every node/.style={scale=#1}]
  \draw (0,0) rectangle (2.4,2.4);
  \node[anchor=north west] at (0.15,2.25) {Area:};
  \node[anchor=east]  at (-0.15,1.2) {#2};
  \node[anchor=north] at (1.2,-0.15) {#2};
  \begin{scope}[xshift=4.6cm]
    \draw (0,0) rectangle (2.4,2.4);
    \node[anchor=north west] at (0.15,2.25) {Area:};
    \node[anchor=east]  at (-0.15,1.2) {#3};
    \node[anchor=north] at (1.2,-0.15) {#3};
  \end{scope}
\end{tikzpicture}
\end{center}}

\providecommand{\samecubes}[3][1]{%
\begin{center}
{\footnotesize These cubes are the \textbf{same size}.}

\vspace{1.5mm}
\begin{tikzpicture}[scale=#1,every node/.style={scale=#1}]
  \foreach \dx in {0,5.2} {
    \begin{scope}[xshift=\dx cm]
      \draw (0,0) rectangle (2.0,2.0);
      \draw (0,2.0) -- (0.65,2.6) -- (2.65,2.6) -- (2.0,2.0);
      \fill[black!12] (2.0,0) -- (2.65,0.6) -- (2.65,2.6) -- (2.0,2.0) -- cycle;
      \draw (2.0,0) -- (2.65,0.6) -- (2.65,2.6) -- (2.0,2.0);
      \node[anchor=north west] at (0.1,1.9) {Volume:};
    \end{scope}
  }
  \node[anchor=east]  at (-0.15,1.0) {#2};
  \node[anchor=north] at (1.0,-0.15) {#2};
  \node[anchor=west]  at (2.25,0.4) {#2};
  \node[anchor=east]  at (5.05,1.0) {#3};
  \node[anchor=north] at (6.2,-0.15) {#3};
  \node[anchor=west]  at (7.45,0.4) {#3};
\end{tikzpicture}
\end{center}}

%% NOTE: do NOT add \usepackage to individual activity .tex files.
%% When a xourse inlines an activity it gobbles \documentclass and \input, but a bare
%% \usepackage still executes -- after \begin{document} -- and errors out.
%% All shared packages and macros belong here.
%%
%% ximera.cls already loads: amsmath, amssymb, amsthm, cancel, enumitem, graphicx,
%% hyperref, listings, pgfplots, tikz, url, xcolor. No need to re-load those.
, so this file is
%% read twice. That was harmless while it held only \usepackage lines (LaTeX skips a
%% package it has already loaded) but a second \newcommand is a hard error.
%%
%%   \blank[width]        fill-in-the-blank rule (default 2.5cm)
%%   \showwork            "show and explain your work" reminder
%%   \showworkline        ... plus which schedule line was used
%%   \showworkyear        ... plus which year's schedule was used
%%   \schedhead{...}      centred bold caption above a tiered-rate table
%%   \samesquares[scale]{left}{right}   two equal-area squares, labelled
%%   \samecubes[scale]{left}{right}     two equal-volume cubes, labelled
%% ---------------------------------------------------------------

\providecommand{\blank}[1][2.5cm]{\underline{\hspace{#1}}}
\providecommand{\showwork}{\textbf{REMEMBER TO SHOW AND EXPLAIN YOUR WORK.}}
\providecommand{\showworkline}{\textbf{REMEMBER TO SHOW AND EXPLAIN YOUR WORK.
  Explanation should include how and why you know which line of the
  schedule was used to calculate this amount.}}
\providecommand{\showworkyear}{\begin{sloppypar}\textbf{REMEMBER TO SHOW AND
  EXPLAIN YOUR WORK. Explanation should include how and why you know which
  line of the year's tax schedule was used to calculate this tax
  amount.}\end{sloppypar}}
\providecommand{\schedhead}[1]{\begin{center}\textbf{#1}\end{center}}

\providecommand{\samesquares}[3][1]{%
\begin{center}
{\footnotesize These squares are the \textbf{same size}.}

\vspace{1.5mm}
\begin{tikzpicture}[scale=#1,every node/.style={scale=#1}]
  \draw (0,0) rectangle (2.4,2.4);
  \node[anchor=north west] at (0.15,2.25) {Area:};
  \node[anchor=east]  at (-0.15,1.2) {#2};
  \node[anchor=north] at (1.2,-0.15) {#2};
  \begin{scope}[xshift=4.6cm]
    \draw (0,0) rectangle (2.4,2.4);
    \node[anchor=north west] at (0.15,2.25) {Area:};
    \node[anchor=east]  at (-0.15,1.2) {#3};
    \node[anchor=north] at (1.2,-0.15) {#3};
  \end{scope}
\end{tikzpicture}
\end{center}}

\providecommand{\samecubes}[3][1]{%
\begin{center}
{\footnotesize These cubes are the \textbf{same size}.}

\vspace{1.5mm}
\begin{tikzpicture}[scale=#1,every node/.style={scale=#1}]
  \foreach \dx in {0,5.2} {
    \begin{scope}[xshift=\dx cm]
      \draw (0,0) rectangle (2.0,2.0);
      \draw (0,2.0) -- (0.65,2.6) -- (2.65,2.6) -- (2.0,2.0);
      \fill[black!12] (2.0,0) -- (2.65,0.6) -- (2.65,2.6) -- (2.0,2.0) -- cycle;
      \draw (2.0,0) -- (2.65,0.6) -- (2.65,2.6) -- (2.0,2.0);
      \node[anchor=north west] at (0.1,1.9) {Volume:};
    \end{scope}
  }
  \node[anchor=east]  at (-0.15,1.0) {#2};
  \node[anchor=north] at (1.0,-0.15) {#2};
  \node[anchor=west]  at (2.25,0.4) {#2};
  \node[anchor=east]  at (5.05,1.0) {#3};
  \node[anchor=north] at (6.2,-0.15) {#3};
  \node[anchor=west]  at (7.45,0.4) {#3};
\end{tikzpicture}
\end{center}}

%% NOTE: do NOT add \usepackage to individual activity .tex files.
%% When a xourse inlines an activity it gobbles \documentclass and \input, but a bare
%% \usepackage still executes -- after \begin{document} -- and errors out.
%% All shared packages and macros belong here.
%%
%% ximera.cls already loads: amsmath, amssymb, amsthm, cancel, enumitem, graphicx,
%% hyperref, listings, pgfplots, tikz, url, xcolor. No need to re-load those.
, so this file is
%% read twice. That was harmless while it held only \usepackage lines (LaTeX skips a
%% package it has already loaded) but a second \newcommand is a hard error.
%%
%%   \blank[width]        fill-in-the-blank rule (default 2.5cm)
%%   \showwork            "show and explain your work" reminder
%%   \showworkline        ... plus which schedule line was used
%%   \showworkyear        ... plus which year's schedule was used
%%   \schedhead{...}      centred bold caption above a tiered-rate table
%%   \samesquares[scale]{left}{right}   two equal-area squares, labelled
%%   \samecubes[scale]{left}{right}     two equal-volume cubes, labelled
%% ---------------------------------------------------------------

\providecommand{\blank}[1][2.5cm]{\underline{\hspace{#1}}}
\providecommand{\showwork}{\textbf{REMEMBER TO SHOW AND EXPLAIN YOUR WORK.}}
\providecommand{\showworkline}{\textbf{REMEMBER TO SHOW AND EXPLAIN YOUR WORK.
  Explanation should include how and why you know which line of the
  schedule was used to calculate this amount.}}
\providecommand{\showworkyear}{\begin{sloppypar}\textbf{REMEMBER TO SHOW AND
  EXPLAIN YOUR WORK. Explanation should include how and why you know which
  line of the year's tax schedule was used to calculate this tax
  amount.}\end{sloppypar}}
\providecommand{\schedhead}[1]{\begin{center}\textbf{#1}\end{center}}

\providecommand{\samesquares}[3][1]{%
\begin{center}
{\footnotesize These squares are the \textbf{same size}.}

\vspace{1.5mm}
\begin{tikzpicture}[scale=#1,every node/.style={scale=#1}]
  \draw (0,0) rectangle (2.4,2.4);
  \node[anchor=north west] at (0.15,2.25) {Area:};
  \node[anchor=east]  at (-0.15,1.2) {#2};
  \node[anchor=north] at (1.2,-0.15) {#2};
  \begin{scope}[xshift=4.6cm]
    \draw (0,0) rectangle (2.4,2.4);
    \node[anchor=north west] at (0.15,2.25) {Area:};
    \node[anchor=east]  at (-0.15,1.2) {#3};
    \node[anchor=north] at (1.2,-0.15) {#3};
  \end{scope}
\end{tikzpicture}
\end{center}}

\providecommand{\samecubes}[3][1]{%
\begin{center}
{\footnotesize These cubes are the \textbf{same size}.}

\vspace{1.5mm}
\begin{tikzpicture}[scale=#1,every node/.style={scale=#1}]
  \foreach \dx in {0,5.2} {
    \begin{scope}[xshift=\dx cm]
      \draw (0,0) rectangle (2.0,2.0);
      \draw (0,2.0) -- (0.65,2.6) -- (2.65,2.6) -- (2.0,2.0);
      \fill[black!12] (2.0,0) -- (2.65,0.6) -- (2.65,2.6) -- (2.0,2.0) -- cycle;
      \draw (2.0,0) -- (2.65,0.6) -- (2.65,2.6) -- (2.0,2.0);
      \node[anchor=north west] at (0.1,1.9) {Volume:};
    \end{scope}
  }
  \node[anchor=east]  at (-0.15,1.0) {#2};
  \node[anchor=north] at (1.0,-0.15) {#2};
  \node[anchor=west]  at (2.25,0.4) {#2};
  \node[anchor=east]  at (5.05,1.0) {#3};
  \node[anchor=north] at (6.2,-0.15) {#3};
  \node[anchor=west]  at (7.45,0.4) {#3};
\end{tikzpicture}
\end{center}}

%% NOTE: do NOT add \usepackage to individual activity .tex files.
%% When a xourse inlines an activity it gobbles \documentclass and \input, but a bare
%% \usepackage still executes -- after \begin{document} -- and errors out.
%% All shared packages and macros belong here.
%%
%% ximera.cls already loads: amsmath, amssymb, amsthm, cancel, enumitem, graphicx,
%% hyperref, listings, pgfplots, tikz, url, xcolor. No need to re-load those.

\title{In-Class: Credit Cards and Mortgages}
\begin{document}
\begin{abstract}
How a credit-card balance grows when you pay only the minimum, why the balance falls
by less than the payment, and what a down payment does to the total cost of a
30-year mortgage.
\end{abstract}
\maketitle

\section*{Credit Cards}

Let's investigate what might happen if we were to not pay off the whole balance on a
credit card by the Payment Due Date.

According to \href{http://CreditCards.com}{CreditCards.com}'s credit card report, the
current national average APR for credit cards is $17.67\%$. Let's assume Little Jimmy
has a credit card with that APR and spent \$850 in his first month having the card.

\begin{problem}
What does APR stand for, and what is the monthly rate for this credit card?

The monthly rate is $\answer[tolerance=0.005]{1.47}$\% per month.

\begin{solution}
APR stands for \textbf{Annual} Percentage Rate --- it is a yearly rate, so to get the
monthly rate we divide by 12:
\[
\frac{\text{APR}}{12 \text{ months}} = \frac{17.67\%}{12 \text{ months}}
= \frac{1.47\%}{\text{month}}.
\]
\end{solution}
\end{problem}

Suppose Jimmy's minimum payment is \$30, and he makes only that minimum payment each
month.

\begin{problem}
Fill in the table below. Each month, interest is charged on the balance left after
the previous payment; that interest is added to the balance, and then the payment is
subtracted.

\begin{center}
\begin{tabular}{|c|l|l|c|l|}
\hline
After \_\_ months & Interest charged & Balance & Payment & Remaining balance \\
\hline
0 & N/A & \$850 & \$30 & $\$850 - \$30 = \$820$ \\
\hline
1 & $\$820(0.0147) = \$12.05$ & $\$832.05$ & \$30 & $\$\answer[tolerance=0.02]{802.05}$ \\
\hline
2 & $\$\answer[tolerance=0.02]{11.79}$ & $\$\answer[tolerance=0.02]{813.84}$ & \$30 & $\$\answer[tolerance=0.02]{783.84}$ \\
\hline
3 & $\$\answer[tolerance=0.02]{11.52}$ & $\$\answer[tolerance=0.02]{795.36}$ & \$30 & $\$\answer[tolerance=0.02]{765.36}$ \\
\hline
\end{tabular}
\end{center}

\begin{hint}
For month 1: the balance after the month-0 payment was \$820. Interest is
$\$820 \cdot 0.0147 = \$12.05$, so the balance becomes
$\$820 + \$12.05 = \$832.05$, and after the \$30 payment
$\$832.05 - \$30 = \$802.05$.
\end{hint}

\begin{solution}
\[
\begin{aligned}
\text{Month 2:} \quad & \$802.05 \cdot 0.0147 = \$11.79, &&
  \$802.05 + \$11.79 = \$813.84, && \$813.84 - \$30 = \$783.84 \\
\text{Month 3:} \quad & \$783.84 \cdot 0.0147 = \$11.52, &&
  \$783.84 + \$11.52 = \$795.36, && \$795.36 - \$30 = \$765.36
\end{aligned}
\]
\end{solution}
\end{problem}

\begin{problem}
Even though Jimmy is paying \$30 each month, how much has the balance actually gone
down each month? Why is it not going down by \$30?

Month 1: $\$\answer[tolerance=0.02]{17.95}$ \qquad
Month 2: $\$\answer[tolerance=0.02]{18.21}$ \qquad
Month 3: $\$\answer[tolerance=0.02]{18.48}$

\begin{freeResponse}
\end{freeResponse}

\begin{solution}
Comparing the balances month to month:
\[
\begin{array}{ll}
\$850 - \$832.05 = \$17.95 & \text{1st month} \\
\$832.05 - \$813.84 = \$18.21 & \text{2nd month} \\
\$813.84 - \$795.36 = \$18.48 & \text{3rd month}
\end{array}
\]
The balance does not go down by the full \$30 because Jimmy is accruing interest.
Part of each payment goes to the interest that was just added, and only what is left
reduces the balance. Notice the amount of progress creeps up each month --- as the
balance shrinks, so does the interest charged on it.
\end{solution}
\end{problem}

We can see from just a few months of calculations how this gets complicated quickly
--- and that is without Jimmy continuing to use the card or varying his payments.
Luckily there are online calculators that can help.

\begin{problem}
Go to the
\href{https://www.bankrate.com/calculators/credit-cards/credit-card-payoff-calculator.aspx}{Bankrate
credit card payoff calculator} and enter Jimmy's situation.

Two notes before you do. The online calculator starts charging interest on the card
balance immediately --- it assumes the balance is the amount \emph{after} the first
payment due date has passed. So you will not be entering \$850 as the card balance.
Also, the calculator wants the APR, not the monthly rate.

What should you enter as the card balance?
\[
\$\answer{820}
\]

\begin{solution}
Enter $\$850 - \$30 = \$820$, because Jimmy does not begin accruing interest until
after the first payment.
\end{solution}
\end{problem}

\begin{problem}
Using the calculator:

\begin{prompt}
\textbf{(a)} How long will it take Jimmy to pay off this debt making only the minimum
\$30 payment?
\begin{freeResponse}
\end{freeResponse}
\end{prompt}

\begin{prompt}
\textbf{(b)} What total dollar amount of interest will he have paid?
\begin{freeResponse}
\end{freeResponse}
\end{prompt}

\begin{prompt}
\textbf{(c)} Suppose Jimmy doubles the minimum payment each month. How long will it
take him to pay off the debt, and what will he pay in total interest?
\begin{freeResponse}
\end{freeResponse}
\end{prompt}

\begin{prompt}
\textbf{(d)} Did doubling the payment cut the payoff time in half? Did it cut the
interest in half?
\begin{freeResponse}
\end{freeResponse}
\end{prompt}

\begin{prompt}
\textbf{(e)} How much must Jimmy pay each month to pay it off in 6 months? In 1 year?
\begin{freeResponse}
\end{freeResponse}
\end{prompt}

\begin{prompt}
\textbf{(f)} Did the monthly amount get cut in half when the time doubled? Explain
how part (d) could have told you in advance whether it would.
\begin{freeResponse}
\end{freeResponse}
\end{prompt}
\end{problem}

\section*{Mortgages}

In the pre-class work, some general advice was given: that no more than $30\%$ to
$40\%$ of one's income should go towards housing payments, either as rent or as a
mortgage. If you read more about these rules, you'll find it's not necessarily
\emph{housing} that should stay under $30\%$--$40\%$ of income, but \emph{total debt}.
There are helpful calculators such as the
\href{https://www.bankrate.com/calculators/mortgages/new-house-calculator.aspx}{Bankrate
new house calculator}, where you can enter all sorts of debt factors alongside a
mortgage.

\begin{problem}
Try inputting Susie and her significant other's combined income of \$103,500 into the
calculator above. Leave the auto-filled values in the ``New Home Info'' section, and
set everything in ``Monthly Expenses'' to \$0.

\begin{prompt}
\textbf{(a)} You should \emph{not} have entered ``\$103,500'' in the wages section.
Why not? What value should you enter?
\begin{freeResponse}
\end{freeResponse}
\end{prompt}

\begin{prompt}
\textbf{(b)} What does the calculator give for an ``Available Mortgage Payment,'' and
what percent of the monthly income is that?
\begin{freeResponse}
\end{freeResponse}
\end{prompt}

\begin{prompt}
\textbf{(c)} Now add values to the ``Monthly Expenses'' section --- a car payment,
credit card payments, maybe student loans under ``other debts.'' Keep everything else
the same. How does this change the Available Mortgage Payment? Record some values
that \emph{do} change it, and some that do \emph{not}.
\begin{freeResponse}
\end{freeResponse}
\end{prompt}
\end{problem}

Now let's focus on the mortgage loan itself. Suppose a buyer is looking to purchase a
home for \$175,000, using the
\href{https://www.bankrate.com/calculators/mortgages/mortgage-calculator-beta/}{Bankrate
mortgage calculator}. Set Homeowner's Insurance, Property Taxes, and HOA fees to \$0
so we can focus on the mortgage itself --- though those values give you an idea of
what else to consider when home-buying. Use a mortgage rate of $4.8\%$ and a 30 year
loan.

\begin{problem}
Suppose the buyer wants to put $20\%$ of the cost down. (Putting $20\%$ or more down
usually means the buyer does not have to pay mortgage insurance on the loan.)

\begin{prompt}
\textbf{(a)} How much would the down payment be?
\[
\$\answer{35000}
\]
\begin{solution}
$20\%$ of \$175,000 is $0.20 \cdot \$175{,}000 = \$35{,}000$.
\end{solution}
\end{prompt}

\begin{prompt}
\textbf{(b)} Put this into the calculator. The monthly mortgage payment is
\[
\$\answer[tolerance=1]{734} \text{ per month.}
\]
\end{prompt}

\begin{prompt}
\textbf{(c)} After the 30 years of the loan, how much in total would the buyer have
paid for the home? Don't forget to include the down payment.
\[
\$\answer[tolerance=10]{299240}
\]

\begin{hint}
First convert 30 years into months.
\end{hint}

\begin{solution}
\[
30 \text{ years} \cdot \frac{12 \text{ months}}{1 \text{ year}} = 360 \text{ months},
\]
\[
\frac{\$734}{\text{month}} \cdot 360 \text{ months} = \$264{,}240,
\]
and adding the down payment,
\[
\$264{,}240 + \$35{,}000 = \$299{,}240.
\]
\end{solution}
\end{prompt}
\end{problem}

\begin{problem}
$20\%$ is a lot for a down payment. Suppose the buyer is able to put together
\$20,000 instead.

\begin{prompt}
\textbf{(a)} What percentage of the home price is this down payment?
\begin{center}
$\answer[tolerance=0.05]{11.43}$\%
\end{center}
\begin{solution}
\$20,000 is what percent of \$175,000? From $\$20{,}000 = x \cdot \$175{,}000$,
\[
x = \frac{\$20{,}000}{\$175{,}000} = 0.1143,
\]
so \$20,000 is $11.43\%$ of \$175,000.
\end{solution}
\end{prompt}

\begin{prompt}
\textbf{(b)} Put \$20,000 into the calculator as the down payment. The monthly
mortgage payment is
\[
\$\answer[tolerance=1]{813} \text{ per month.}
\]
\end{prompt}

\begin{prompt}
\textbf{(c)} After 30 years, how much in total would the buyer have paid, including
the down payment?
\[
\$\answer[tolerance=10]{312680}
\]
\begin{solution}
\[
\frac{\$813}{\text{month}} \cdot 360 \text{ months} = \$292{,}680,
\]
\[
\$292{,}680 + \$20{,}000 = \$312{,}680.
\]
\end{solution}
\end{prompt}

\begin{prompt}
\textbf{(d)} How much more will the buyer pay over the 30 year mortgage by making the
smaller down payment?
\[
\$\answer[tolerance=10]{13440}
\]
\begin{solution}
\[
\begin{array}{ll}
\$35{,}000 \text{ down payment} & \text{Total: } \$299{,}240 \\
\$20{,}000 \text{ down payment} & \text{Total: } \$312{,}680
\end{array}
\]
\[
\$312{,}680 - \$299{,}240 = \$13{,}440.
\]
So putting down \$20,000 instead of \$35,000 in this scenario costs an additional
\$13,440 overall --- notice that is nearly as much as the \$15,000 the buyer
``saved'' up front.
\end{solution}
\end{prompt}
\end{problem}

\end{document}
